% =====================================================================
% Slide deck for 5--6 min talk on the closed-loop game-design
% optimisation project. Plain Beamer (no theme), 16:9.
%
% Compile:  pdflatex slides.tex   (run twice if you add \ref/\cite)
% =====================================================================
\documentclass[10pt,aspectratio=169]{beamer}

\usepackage{graphicx}
\usepackage{amsmath, amssymb}
\usepackage{booktabs}
\usepackage{hyperref}
\hypersetup{colorlinks=true,linkcolor=black,urlcolor=blue}

% Plain look: black on white, no nav decoration, no header/footer chrome.
\setbeamertemplate{navigation symbols}{}
\setbeamercolor{normal text}{fg=black,bg=white}
\setbeamercolor{frametitle}{fg=black,bg=white}
\setbeamercolor{title}{fg=black,bg=white}
\setbeamercolor{author}{fg=black}
\setbeamercolor{institute}{fg=black}
\setbeamercolor{date}{fg=black}
\setbeamercolor{itemize item}{fg=black}
\setbeamercolor{itemize subitem}{fg=black}
\setbeamerfont{frametitle}{size=\Large,series=\bfseries}
\setbeamertemplate{itemize item}{$\bullet$}
\setbeamertemplate{itemize subitem}{$\circ$}
\setbeamertemplate{footline}{}  % no footer

\graphicspath{%
  {figures/}%
  {figures/llm/}%
  {figures/boards_mask/}%
  {../final_report/figures/}%
}

\title{Beta-Testing Monopoly}
\subtitle{Closed-Loop Game-Design Optimisation over a Diverse Strategy Pool}
\author{Selim Emir Can}
\institute{Stanford CS348K}
\date{June 2026}

\begin{document}

% ===================== 1. Title =====================
\begin{frame}
  \titlepage
  \begin{center}
    \small \url{github.com/Selim-Emir-Can/CS348K-proj}
  \end{center}
\end{frame}

% ===================== 2. Motivation =====================
\begin{frame}{Motivation}
  \begin{columns}[T,onlytextwidth]
    \begin{column}{0.55\textwidth}
      \begin{itemize}
        \item Game design is a \textbf{parameter-rich search problem}: every cost, reward, payoff, and structural rule is a designer knob.
        \item Human playtesting can't span that space, so automated agents have served as designer-side proxies for over a decade: procedural content generation, search-based game design, RL playtesters, generative agents, league play (AlphaStar).
        \item These pipelines almost always evaluate against \textbf{one strong agent}: one rating, one set of preferences, one blind spot.
        \item \textbf{Question:} can a \emph{diverse pool} of imperfect agents be a more honest design probe than any single oracle?
      \end{itemize}
    \end{column}
    \begin{column}{0.4\textwidth}
      \centering
      \includegraphics[width=\linewidth]{board_default.png}\\
      \tiny Monopoly as our concrete substrate
    \end{column}
  \end{columns}
\end{frame}

% ===================== 3. Problem =====================
\begin{frame}{Problem}
  \begin{itemize}
    \item \textbf{Find boards that are \emph{fair, decisive, paced, interactive}} --- robustly across many player strategies, not just one. Evaluated at 2p and 3p.
    \item \textbf{5-axis composite objective:} fairness, worst-pair fairness, length, draw rate, money transfer. Every weight and target exposed as a CLI flag, so ablations are one-line runs.
    \item \textbf{Where prior automated playtesting falls short:} most pipelines use a single RL agent or LLM as evaluator --- one rating, one set of preferences, no way to detect blind spots. \emph{Cross-class agreement} is the missing piece.
  \end{itemize}
\end{frame}

% ===================== 4. Our approach =====================
\begin{frame}{Our approach}
  \begin{columns}[T,onlytextwidth]
    \begin{column}{0.55\textwidth}
      \begin{itemize}
        \item \textbf{Board space} (66 dims): per-property cost $\times$ rent multipliers + 22-bit keep-mask (lap shrinkage).
        \item \textbf{Strategy pool} (30): 10 hand-designed archetypes + 20 parametric samples, 17-dim player parameterisation.
        \item \textbf{Two evaluator classes:} rule-based pool, and Qwen2.5-1.5B with deterministic STATE/ECHO/REASON/ANSWER validation.
        \item \textbf{GA} (pop $30 \times$ gens $30$) with CRN seed discipline: same (vec, seeds) $\rightarrow$ byte-identical metrics.
      \end{itemize}
    \end{column}
    \begin{column}{0.4\textwidth}
      \centering
      \includegraphics[width=\linewidth]{board_ga_2p_mask.png}\\
      \tiny default vs.\ GA-2p winning board
    \end{column}
  \end{columns}
\end{frame}

% ===================== 5. The strategy pool =====================
\begin{frame}{The strategy pool: 10 named archetypes (+ 20 sampled)}
  \centering
  \scriptsize
  \begin{tabular}{lccccl}
    \toprule
    Strategy & Cash floor & Trades & Build & Buys & Ignores groups \\
    \midrule
    AggressiveBuilder & $0$    & no  & max  & RR+Util & --- \\
    CashHoarder       & $1000$ & no  & slow & RR+Util & --- \\
    Trader            & $300$  & yes & max  & RR+Util & --- \\
    RailroadKing      & $100$  & no  & max  & RR only & all colour groups \\
    LowCostOnly       & $200$  & yes & max  & RR+Util & 5 priciest (Or/Rd/Yl/Gn/In) \\
    HighCostOnly      & $400$  & yes & max  & RR only & 3 cheapest (Br/Lb/Pk) \\
    Balanced          & $300$  & yes & max  & RR+Util & --- \\
    Passive           & $500$  & no  & slow & none    & --- \\
    Bully             & $50$   & yes & max  & RR+Util & --- \\
    RiskAverse        & $800$  & no  & slow & RR only & Green, Indigo \\
    \bottomrule
  \end{tabular}
  \vspace{0.7em}
  \begin{itemize}\small
    \item Plus \textbf{20 sampled} strategies drawn from the 17-dim space (Bernoulli booleans, random cash/trade thresholds, each colour group ignored w.p.\ $0.25$).
    \item ``Cash floor'' $=$ cash kept un-spent; ``Build'' $=$ max houses/turn vs.\ one/turn; ``Buys'' $=$ railroads / utilities.
  \end{itemize}
\end{frame}

% ===================== 6. Rule-based GA: convergence & Pareto =====================
\begin{frame}{Rule-based GA: convergence \& Pareto trade-off}
  \begin{columns}[T,onlytextwidth]
    \begin{column}{0.38\textwidth}
      \begin{itemize}\small
        \item \textbf{$+47\%$} over default at both 2p and 3p.
        \item \textbf{$+13\%$} over random search at matched $\sim$842-eval budget.
        \item Pareto: \emph{fairness vs.\ rounds} is the dominant trade-off; ablations populate different regions of the front.
        \item 2p$\leftrightarrow$3p cross-eval: $16$--$24\%$ off-regime gap; both designs still beat default.
      \end{itemize}
    \end{column}
    \begin{column}{0.6\textwidth}
      \centering
      \begin{tabular}{cc}
        \includegraphics[width=0.46\linewidth]{convergence_2p.png} &
        \includegraphics[width=0.46\linewidth]{convergence_3p.png} \\
        \tiny convergence, 2p & \tiny convergence, 3p \\[0.2em]
        \includegraphics[width=0.46\linewidth]{ga_2p_pareto.png} &
        \includegraphics[width=0.46\linewidth]{ga_3p_pareto.png} \\
        \tiny Pareto front, 2p & \tiny Pareto front, 3p \\
      \end{tabular}
    \end{column}
  \end{columns}
\end{frame}

% ===================== 7. What the GA learned about the design space =====================
\begin{frame}{What the GA learned about Monopoly}
  \begin{columns}[T,onlytextwidth]
    \begin{column}{0.38\textwidth}
      \begin{itemize}\small
        \item \textbf{Cost \& rent multipliers:} GA inflates high-value colour groups, deflates low-value ones --- ``rich-get-richer'' is dampened, not eliminated.
        \item \textbf{Keep-mask:} shortens the lap by removing the cheapest properties; each lap becomes more consequential.
        \item \textbf{$30\!\times\!30$ win-rate diff:} which strategy matchups got fairer (blue) vs.\ stayed locked (red). Bulk of cells shift toward fairness; a few are immutable.
      \end{itemize}
    \end{column}
    \begin{column}{0.6\textwidth}
      \centering
      \begin{tabular}{cc}
        \includegraphics[width=0.46\linewidth]{cost_multipliers_2p.png} &
        \includegraphics[width=0.46\linewidth]{rent_multipliers_2p.png} \\
        \tiny cost multipliers, 2p & \tiny rent multipliers, 2p \\[0.2em]
        \includegraphics[width=0.46\linewidth]{keep_mask_2p.png} &
        \includegraphics[width=0.46\linewidth]{heatmap_ga2p.diff.png} \\
        \tiny keep mask, 2p & \tiny win-rate diff (GA $-$ default) \\
      \end{tabular}
    \end{column}
  \end{columns}
\end{frame}

% ===================== 8. LLM cross-class probe =====================
\begin{frame}{LLM cross-class probe: directional agreement}
  \begin{columns}[T,onlytextwidth]
    \begin{column}{0.4\textwidth}
      \begin{itemize}\small
        \item Qwen2.5-1.5B with all-LLM seats re-evaluates the GA winners as an \textbf{independent agent class}.
        \item Reproduces the rule-based directional ranking: rounds $-50\%$ at 2p on the GA-winner; transfer increases at 3p.
        \item Two independent agent classes agree $\Rightarrow$ the GA winner is \textbf{not an artifact of the evaluator}.
        \item Asymmetry: LLM tracks the design improvement more strongly in one regime than the other.
      \end{itemize}
    \end{column}
    \begin{column}{0.58\textwidth}
      \centering
      \includegraphics[width=0.85\linewidth]{fig_cross_class_agreement.png}\\
      \tiny directional agreement (rule-based vs.\ LLM)
      \vspace{0.4em}\\
      \includegraphics[width=0.85\linewidth]{fig_fairness_asymmetry.png}\\
      \tiny fairness asymmetry across regimes
    \end{column}
  \end{columns}
\end{frame}

% ===================== 9. LLM-driven GA =====================
\begin{frame}{LLM-driven GA: when the LLM \emph{is} the evaluator}
  \begin{columns}[T,onlytextwidth]
    \begin{column}{0.38\textwidth}
      \begin{itemize}\small
        \item Swapped the evaluator class to all-LLM seats and \textbf{re-ran the GA} (pop $8 \times$ gens $5$).
        \item The LLM-driven loop converges cleanly to a board the \textbf{LLM rates highly}.
        \item Score distribution shows the LLM is internally consistent about what it prefers.
        \item Question: does that board hold up under a \emph{different} evaluator?
      \end{itemize}
    \end{column}
    \begin{column}{0.6\textwidth}
      \centering
      \begin{tabular}{cc}
        \includegraphics[width=0.46\linewidth]{fig_llm_ga_convergence.png} &
        \includegraphics[width=0.46\linewidth]{fig_llm_ga_score_distribution.png} \\
        \tiny LLM-driven GA convergence & \tiny score distribution \\[0.3em]
        \multicolumn{2}{c}{\includegraphics[width=0.55\linewidth]{board_llm_ga_2p.png}} \\
        \multicolumn{2}{c}{\tiny board the LLM-driven GA converged to} \\
      \end{tabular}
    \end{column}
  \end{columns}
\end{frame}

% ===================== 10. Why single-oracle evaluators fall short =====================
\begin{frame}{Why single-oracle evaluators fall short}
  \begin{columns}[T,onlytextwidth]
    \begin{column}{0.48\textwidth}
      \begin{itemize}
        \item Cross-evaluate that LLM-driven GA winner with the 30-strategy rule-based pool.
        \item \textbf{Fairness gap:} $0.20$ (LLM rating) vs.\ $0.379$ (pool) at 2p; $0.05$ vs.\ $0.404$ at 3p.
        \item Gap \textit{widens} with player count: a single-personality evaluator misses more as the space gets richer.
        \item \textbf{The board the LLM loved is not actually fair} --- you only see it with a second agent class.
      \end{itemize}
    \end{column}
    \begin{column}{0.5\textwidth}
      \centering
      \includegraphics[width=\linewidth]{fig_cross_evaluator_gap.png}
    \end{column}
  \end{columns}
\end{frame}

% ===================== 11. Takeaway =====================
\begin{frame}{Takeaway}
  \begin{itemize}
    \item \textbf{Pool of strategies $>$ single oracle} as an evaluator for any multi-agent design problem.
    \item Cross-class probing \textbf{quantifies blind spots}: you can't see what your evaluator misses if it's the only one looking.
    \item Variance-reduction discipline (CRN) makes small-budget GA see real signal; the bottleneck is problem framing, not optimiser cleverness.
    \item Transfers to market mechanisms, recommendation systems, RL environment design.
  \end{itemize}
  \vspace{1.2em}
  \begin{center}
    \small Code, logs, figures: \url{github.com/Selim-Emir-Can/CS348K-proj}
  \end{center}
\end{frame}

\end{document}
